\documentclass[ ../main.tex]{subfiles}
\providecommand{\mainx}{..}
\begin{document}
\section{Perfect hash functions}
A \emph{\gls{gls-set}} is an unordered collection of distinct elements.
If we know the elements in a set, we may denote the set by these elements, e.g., $\{a,c,b\}$ denotes a set whose membes are exactly $a$, $b$, and $c$.

A \emph{\gls{gls-finiteset}} has a finite number of elements.
For example, $\{ 1, 3, 5 \}$ is a finite set with three elements.
When sets $\Set{A}$ and $\Set{B}$ are \emph{isomorphic}, denoted by $\Set{A} 
\cong \Set{B}$, they can be put into a one-to-one correspondence (bijection), e.g., $\{b,a,c\} \cong \{1,2,3\}$.
Since there exists at least one bijection between isomorphic sets, we can losslessly convert one to the other and thus, isomorphic sets are in some sense equivalent.

The cardinality of a finite set $\Set{A}$ is the number of elements in the set, denoted by $\Card{\Set{A}}$, e.g., $\Card{\left\{ 1, 3, 5\right\}} = 3$.
A \emph{countably infinite set} is isomorphic to the set of \emph{natural numbers} $\NatSet \coloneqq \{1,2,3,4,5,\ldots\}$.

Given two elements $a$ and $b$, an ordered pair of $a$ then $b$ is denoted by $\Pair{a}{b}$, where $\Pair{a}{b} = \Pair{c}{d}$ \emph{if and only if} $a = c$ and $b = d$.
Ordered pairs are non-commutative and non-associative, i.e., $\Pair{a}{b} \neq \Pair{b}{a}$ if $a \neq b$ and $\Pair{a}{\Pair{b}{c}} \neq \Pair{\Pair{b}{a}}{c}$.

Related to the ordered pair is the Cartesian product.
\begin{definition}
The set $\Set{X} \times \Set{Y} \coloneqq \left\{\Pair{x}{y} \colon x \in 
\Set{X} \land y \in \Set{Y}\right\}$ is the Cartesian product of sets $\Set{X}$ 
and $\Set{Y}$.
\end{definition}
By the non-commutative and non-associative property of ordered pairs, the Cartesian product is non-commutative and non-associative.
However, they are isomorphic, i.e., $\Set{X} \times \Set{Y} \cong \Set{Y} \times \Set{X}$.

A \emph{tuple} is a generalization of order pairs which can consist of an arbitrary number of elements, e.g., $\Tuple{x_1,x_2,\ldots,x_n}$.
\begin{definition}[$n$-fold Cartesian product]
The $n$-ary Cartesian product of sets $\Set{X}[1],\ldots,\Set{X}[n]$, is given by $\Set{X}[1] \times \cdots \times \Set{X}[n] = \left\{\Tuple{x_1,\ldots,x_n} \colon x_1 \in \Set{X}[1] \land \cdots \land x_n \in \Set{X}[n]\right\}$.
\end{definition}
Note that $\Set{X}[1] \times \Set{X}[2] \times \Set{X}[3] \cong \Set{X}[1] \times \left(\Set{X}[2] \times \Set{X}[3]\right) \cong \left(\Set{X}[1] \times \Set{X}[2]\right) \times \Set{X}[3]$, thus we may implicitly convert between them without ambiguity.

If each set in the $n$-ary Cartesian product is the same, the power notation may be used, e.g., $\Set{X}^3 \equiv \Set{X} \times \Set{X} \times \Set{X}$.
As special cases, the $0$-ary (nullary) Cartesian product is defined to be $\{ \EmptySet \}$, and the $1$-ary (unary) Cartesian product is the identity, e.g., $\Set{X}^1 = \Set{X}$.



%\begin{example}
%Consider the number $100$ (base $10$). To determine the bit string $100'$, we perform the following steps:
%\begin{enumerate}
%    \item The length of $100'$ is $\lfloor \log_2 100 \rfloor = 6$.
%    \item Subtract $2^6 = 64$ from $100$ which results in $36$.
%    \item $36_{10}$ in binary is $100100_2$.
%    \item $100' = (1\,0\,0\,1\,0\,0)$.
%\end{enumerate}
%\end{example}

%A similar mapping may be defined for $\BitSet^m$ and $\{1,2,\ldots,2^m\}$ as given by
%\begin{equation}
%    b_m \mapsto \sum_{j=1}^{m} 2^{m - j + 1} b_j\,,
%\end{equation}
%where $b_m \in \BitSet^m$.
%   1   empty
%   2   0       
%   3   1
%   4   00
%   5   01
%   6   10
%   7   11
%   8   000
%   9   001
%   10  010
%   11  011
%   12  100
%   13  101
%   14  110
%   15  111


The \emph{\gls{gls-hashfunction}} is given by the following definition.
\begin{definition}
\emph{Hash functions} of type $\Set{X} \mapsto \Set{Y}$ are just total functions, normally with a finite codomain, with the weak assumption that they will be used as a device to assign elements from $\Set{X}$ a value from $\Set{Y}$ without requiring any particular rule.
\end{definition}
For a given bit string $x$ and hash function $\hash$, $y = \hash(x)$ is denoted the \emph{hash} of $x$.

We are particularly interested in \emph{perfect} hash functions as given by the following definition.
\begin{definition}[Perfect hash function]
A \emph{\gls{gls-ph}} over the set $\Set{A} \subseteq \Set{X}$, denoted by
\begin{equation}
    \hash[\Set{A}] \colon \Set{X} \mapsto \Set{Y}\,,
\end{equation}
is an injective function when restricted to $\Set{A}$.\footnote{There are no collisions among elements of $\Set{A}$, $\hash[\Set{A}](x) \neq \hash[\Set{A}](y)$ for all $x, y \in \Set{A}$, $x \neq y$.}
\end{definition}

\begin{assumption}
Perfect hash functions are \emph{surjective}.
\end{assumption}

The \emph{load factor} is given by the following definition.
\begin{definition}
\label{def:loadfactor}
The proportion of hashes mapped to by $\hash_{\Set{A}} \colon \Set{X} \mapsto \NatSet$ over subset $\Set{A}$ is denoted the \gls{gls-loadfactor}.
Specifically, the load factor of $\hash_{\Set{A}}$ is a rational number given by
\begin{equation}
\label{eq:loadfactor}
    \frac{\Card{\hash_{\Set{A}}(\Set{A})}}{\Card{\Image(\hash_{\Set{A}})}}\,.
\end{equation}
\end{definition}


\begin{notation}
A perfect hash function of type $\Set{X} \mapsto \Set{Y}$ over $\Set{A}$ with a load factor $r$ may be denoted by $\hash_{\Set{A}}^{r}$.
If $m = \Card{\Set{A}}$ and we are interested in drawing attention to the cardinality of the perfect hash function, we may also denote it by $\hash_{m}^{r}$.
\end{notation}





\begin{example}
Consider the set $\Set{X} \coloneqq \{x_1,x_2,x_3\}$.
A \gls{gls-ph} of type $\Set{U} \mapsto \NatSet$ over $\Set{X}$ with a load factor $r \coloneqq \frac{3}{5}$ is denoted by $\hash_{\Set{X}}^{0.6}$ or $\hash_{3}^{0.6}$.
Given the load factor $r$ and $\Set{X}$, we may deduce the codomain of $\hash_{\Set{X}}^{r}$ to be precisely $\{0,1,2,3,4\}$.
See \cref{fig:perfhash} for an illustration.
\end{example}

\begin{definition}
A \emph{minimal} perfect hash has a load factor $1$.
\end{definition}

\begin{figure}
\caption{A perfect hash function $\hash_{\Set{X}}^{0.6} \colon \Set{U} \mapsto \{0,1,2,3,4\}$.}
\label{fig:perfhash}
\center
%% Creator: Inkscape inkscape 0.92.1, www.inkscape.org
%% PDF/EPS/PS + LaTeX output extension by Johan Engelen, 2010
%% Accompanies image file 'pf2.pdf' (pdf, eps, ps)
%%
%% To include the image in your LaTeX document, write
%%   \input{<filename>.pdf_tex}
%%  instead of
%%   \includegraphics{<filename>.pdf}
%% To scale the image, write
%%   \def\svgwidth{<desired width>}
%%   \input{<filename>.pdf_tex}
%%  instead of
%%   \includegraphics[width=<desired width>]{<filename>.pdf}
%%
%% Images with a different path to the parent latex file can
%% be accessed with the `import' package (which may need to be
%% installed) using
%%   \usepackage{import}
%% in the preamble, and then including the image with
%%   \import{<path to file>}{<filename>.pdf_tex}
%% Alternatively, one can specify
%%   \graphicspath{{<path to file>/}}
%% 
%% For more information, please see info/svg-inkscape on CTAN:
%%   http://tug.ctan.org/tex-archive/info/svg-inkscape
%%
\begingroup%
  \makeatletter%
  \providecommand\color[2][]{%
    \errmessage{(Inkscape) Color is used for the text in Inkscape, but the package 'color.sty' is not loaded}%
    \renewcommand\color[2][]{}%
  }%
  \providecommand\transparent[1]{%
    \errmessage{(Inkscape) Transparency is used (non-zero) for the text in Inkscape, but the package 'transparent.sty' is not loaded}%
    \renewcommand\transparent[1]{}%
  }%
  \providecommand\rotatebox[2]{#2}%
  \ifx\svgwidth\undefined%
    \setlength{\unitlength}{178.44486146bp}%
    \ifx\svgscale\undefined%
      \relax%
    \else%
      \setlength{\unitlength}{\unitlength * \real{\svgscale}}%
    \fi%
  \else%
    \setlength{\unitlength}{\svgwidth}%
  \fi%
  \global\let\svgwidth\undefined%
  \global\let\svgscale\undefined%
  \makeatother%
  \begin{picture}(1,0.66884337)%
    \put(0.84730911,0.50904772){\color[rgb]{0,0,0}\makebox(0,0)[lb]{\smash{$0$}}}%
    \put(0.84730911,0.3928877){\color[rgb]{0,0,0}\makebox(0,0)[lb]{\smash{$1$}}}%
    \put(0.84730911,0.2767279){\color[rgb]{0,0,0}\makebox(0,0)[lb]{\smash{$2$}}}%
    \put(0.84730911,0.16056788){\color[rgb]{0,0,0}\makebox(0,0)[lb]{\smash{$3$}}}%
    \put(0.84730911,0.04440763){\color[rgb]{0,0,0}\makebox(0,0)[lb]{\smash{$4$}}}%
    \put(0.185,0.631){\color[rgb]{0,0,0}\makebox(0,0)[lb]{\smash{$\Set{X}$}}}%
    \put(0,0){\includegraphics[width=\unitlength,page=1]{pf.pdf}}%
    \put(0.11117223,0.28933608){\color[rgb]{0,0,0}\makebox(0,0)[lb]{\smash{$a$}}}%
    \put(0.17763852,0.4254032){\color[rgb]{0,0,0}\makebox(0,0)[lb]{\smash{$b$}}}%
    \put(0.22523437,0.15989424){\color[rgb]{0,0,0}\makebox(0,0)[lb]{\smash{$c$}}}%
    \put(0.806,0.631){\color[rgb]{0,0,0}\makebox(0,0)[lb]{\smash{hash}}}%
    \put(0,0){\includegraphics[width=\unitlength,page=2]{pf.pdf}}%
    \put(0.433,0.12){\color[rgb]{0,0,0}\makebox(0,0)[lb]{\smash{$\mathsmaller{\hash_{\Set{X}}^{0.6}(c)}$}}}%
    \put(0,0){\includegraphics[width=\unitlength,page=3]{pf.pdf}}%
  \end{picture}%
\endgroup%

\end{figure}

Every hash function in $\Set{X} \mapsto \Set{Y}$ is a perfect hash function over some subset of $\Set{X}$, e.g., every hash function is trivially a perfect hash function of $\EmptySet$ and singleton sets.

Assuming $\Set{X}$ and $\Set{Y}$ are finite, the set of hash functions of type $\Set{X} \mapsto \Set{Y}$, which may also be denoted by $\Set{Y}^{\Set{X}}$, has a cardinality
\begin{equation}
\Card{\Set{Y}}^{\Card{\Set{X}}}\,.
\end{equation}
The set of \emph{perfect} hash functions over $\Set{A} \subseteq \Set{X}$ is a subset of $\Set{X} \mapsto \Set{Y}$ with the predicate that no collisions may occur on any pair of elements in $\Set{A}$.
The set of perfect hash functions over $\Set{A}$ has a cardinality
\begin{equation}
\Perm{\Card{\Set{Y}}}{\Card{\Set{A}}} \Card{\Set{Y}}^{\Card{\Set{X}} - \Card{\Set{A}}}\,.
\end{equation}
\end{document}
